\documentclass[12pt,a4paper]{article}
\usepackage[utf8]{vietnam}
\usepackage{amsmath,amssymb}
\usepackage{tabularx}
\usepackage{tikz}
\usepackage{xcolor}
\usepackage[top=0.6cm, bottom=0.6cm, left=0.6cm, right=0.6cm]{geometry}
\usepackage[dethi]{ex_test}

\SAOPTN{widthSA=3.2,heightSA=0.62}

\setlength{\parskip}{0pt}

% Dòng kẻ chấm mờ đều đặn chuẩn viết tay
\newcommand{\dongke}[1]{%
  \begingroup%
  \setlength{\parskip}{0pt}%
  \baselineskip=18.5pt%
  \vspace{1.0mm}%
  \foreach \i in {1,...,#1}{%
    \noindent{\color{gray!60}\dotfill}\par%
  }%
  \vspace{0.6mm}%
  \endgroup%
}

% Định nghĩa ô bài tập tự luận trong bảng
\newcommand{\oBaiTap}[3][6]{%
  \vspace{1.0mm}%
  \noindent\textbf{#2} #3\par%
  \dongke{#1}%
}

\begin{document}
\pagestyle{empty}
% =============================================================================
% ========================= TOÁN 9 - TUẦN 07 - BUỔI 02: MÃ ĐỀ 101 =============
% ========================= MẶT TRƯỚC (TRANG 1/2) ==============================
% =============================================================================

\noindent
\begin{tabularx}{\linewidth}{|X|c|}
\hline
\textbf{PHIẾU KIỂM TRA ĐỊNH KỲ TOÁN 9 - TUẦN 07 - BUỔI 02} & \textbf{Mã đề: 101 - Trang 1/2} \\
\hline
\end{tabularx}

\vspace{1.0mm}

\noindent
\begin{tabularx}{\linewidth}{|p{3.2cm}|X|}
\hline
\centering\textbf{ĐIỂM SỐ} & \centering\textbf{LỜI NHẬN XÉT CỦA THẦY CÔ} \tabularnewline
\hline
\rule{0pt}{1.8cm} & \tabularnewline
\hline
\end{tabularx}

\vspace{1.2mm}
\noindent\textbf{Câu 1 (0,25 điểm).} $\Delta ABC$ vuông tại $A$ ($\widehat{B}=30^\circ, BC=12\text{ cm}$). Cạnh $AC =$ \shortans{6\text{ cm}}\par\vspace{0.6mm}
\noindent\textbf{Câu 2 (0,25 điểm).} Cho góc nhọn $\alpha$ có $\sin \alpha = \dfrac{3}{5}$. Khi đó $\cos \alpha =$ \shortans{\dfrac{4}{5}}\par\vspace{0.6mm}
\noindent\textbf{Câu 3 (0,25 điểm).} Cho $\Delta ABC$ vuông tại $A$, đường cao $AH$. Biết $BH=4, CH=9$. Khi đó $AH =$ \shortans{6}\par\vspace{0.6mm}
\noindent\textbf{Câu 4 (0,25 điểm).} Rút gọn biểu thức $\sqrt{(2 - \sqrt{5})^2}$ được kết quả: \shortans{\sqrt{5} - 2}\par\vspace{0.6mm}
\noindent\textbf{Câu 5 (0,25 điểm).} Biểu thức $\sqrt{3x - 6}$ xác định khi: \shortans{x \ge 2}\par\vspace{0.6mm}
\noindent\textbf{Câu 6 (0,25 điểm).} Giá trị của biểu thức $\sqrt{18} : \sqrt{2} + \sqrt{16}$ là: \shortans{7}\par\vspace{0.6mm}
\noindent\textbf{Câu 7 (0,25 điểm).} Nghiệm của phương trình $\sqrt{2x - 1} = 3$ là $x =$ \shortans{5}\par\vspace{0.6mm}
\noindent\textbf{Câu 8 (0,25 điểm).} Thang dài $5\text{ m}$ nghiêng $60^\circ$ với đất. Chân thang cách tường: \shortans{2{,}5\text{ m}}\par\vspace{0.6mm}
\noindent\textbf{Câu 9 (0,25 điểm).} Nghiệm của hệ $\begin{cases} x + y = 5 \\ x - y = 1 \end{cases}$ là $(x; y) =$ \shortans{(3; 2)}\par\vspace{0.6mm}
\noindent\textbf{Câu 10 (0,25 điểm).} Giá trị $\sin^2 10^\circ + \dots + \sin^2 80^\circ$ (8 số hạng) là: \shortans{4}\par\vspace{0.6mm}
\vspace{0.8mm}
\noindent\textbf{Câu 11 (2,5 điểm). Đại số tổng hợp:}
\vspace{0.6mm}

\noindent
\begin{tabularx}{\linewidth}{|X|X|}
\hline
\oBaiTap[20]{a) (1,25 điểm)}{$\begin{cases} 2x + 3y = 12 \\ 3x - y = 7 \end{cases}$} & 
\oBaiTap[20]{b) (1,25 điểm)}{Rút gọn với $x > 0, x \ne 4$:\par $P = \left(\dfrac{1}{\sqrt{x}-2} + \dfrac{1}{\sqrt{x}+2}\right) \cdot \dfrac{x-4}{\sqrt{x}}$} \\
\hline
\end{tabularx}

\newpage
% =============================================================================
% ========================= TOÁN 9 - TUẦN 07 - BUỔI 02: MÃ ĐỀ 101 =============
% ========================= MẶT SAU (TRANG 2/2) ================================
% =============================================================================

\noindent
\begin{tabularx}{\linewidth}{|X|c|}
\hline
\textbf{TOÁN 9 - TUẦN 07 - BUỔI 02} & \textbf{Mã đề: 101 - Trang 2/2} \\
\hline
\end{tabularx}

\vspace{1.2mm}

\noindent
\begin{tabularx}{\linewidth}{|X|}
\hline
\vspace{0.5mm}
\textbf{Câu 12 (2,5 điểm).} Hai xe máy cùng khởi hành một lúc từ hai địa điểm $A$ và $B$ cách nhau $90\text{ km}$, đi ngược chiều nhau và gặp nhau sau $1{,}5$ giờ. Biết rằng vận tốc của xe đi từ $A$ lớn hơn vận tốc của xe đi từ $B$ là $10\text{ km/h}$. Tính vận tốc của mỗi xe máy.\par
\vspace{0.4mm}
\noindent\textit{Bài giải:}\dongke{16}
\\ \hline
\end{tabularx}

\vspace{2.0mm}

\noindent
\begin{tabularx}{\linewidth}{|X|}
\hline
\vspace{0.5mm}
\textbf{Câu 13 (2,5 điểm).} Cho tam giác $ABC$ vuông tại $A$ có $\widehat{B} = 30^\circ$ và cạnh huyền $BC = 12\text{ cm}$. Kẻ đường cao $AH$ ($H \in BC$).\par
\textbf{a) (1,5 điểm)} Tính số đo góc $C$, độ dài các cạnh $AC, AB$ và độ dài đường cao $AH$.\par
\textbf{b) (1,0 điểm)} Kẻ $HD \perp AB$ tại $D$ và $HE \perp AC$ tại $E$. Chứng minh rằng: $AD \cdot AB = AE \cdot AC$.\par
\vspace{0.4mm}
\noindent\textit{Bài giải:}\dongke{17}
\\ \hline
\end{tabularx}

\newpage
% =============================================================================
% ========================= TOÁN 9 - TUẦN 07 - BUỔI 02: MÃ ĐỀ 102 =============
% ========================= MẶT TRƯỚC (TRANG 1/2) ==============================
% =============================================================================

\noindent
\begin{tabularx}{\linewidth}{|X|c|}
\hline
\textbf{PHIẾU KIỂM TRA ĐỊNH KỲ TOÁN 9 - TUẦN 07 - BUỔI 02} & \textbf{Mã đề: 102 - Trang 1/2} \\
\hline
\end{tabularx}

\vspace{1.0mm}

\noindent
\begin{tabularx}{\linewidth}{|p{3.2cm}|X|}
\hline
\centering\textbf{ĐIỂM SỐ} & \centering\textbf{LỜI NHẬN XÉT CỦA THẦY CÔ} \tabularnewline
\hline
\rule{0pt}{1.8cm} & \tabularnewline
\hline
\end{tabularx}

\vspace{1.2mm}
\noindent\textbf{Câu 1 (0,25 điểm).} $\Delta ABC$ vuông tại $A$ ($\widehat{B}=30^\circ, BC=16\text{ cm}$). Cạnh $AC =$ \shortans{8\text{ cm}}\par\vspace{0.6mm}
\noindent\textbf{Câu 2 (0,25 điểm).} Cho góc nhọn $\alpha$ có $\cos \alpha = \dfrac{3}{5}$. Khi đó $\sin \alpha =$ \shortans{\dfrac{4}{5}}\par\vspace{0.6mm}
\noindent\textbf{Câu 3 (0,25 điểm).} Cho $\Delta ABC$ vuông tại $A$, đường cao $AH$. Biết $BH=2, CH=8$. Khi đó $AH =$ \shortans{4}\par\vspace{0.6mm}
\noindent\textbf{Câu 4 (0,25 điểm).} Rút gọn biểu thức $\sqrt{(3 - \sqrt{10})^2}$ được kết quả: \shortans{\sqrt{10} - 3}\par\vspace{0.6mm}
\noindent\textbf{Câu 5 (0,25 điểm).} Biểu thức $\sqrt{2x - 8}$ xác định khi: \shortans{x \ge 4}\par\vspace{0.6mm}
\noindent\textbf{Câu 6 (0,25 điểm).} Giá trị của biểu thức $\sqrt{50} : \sqrt{2} - \sqrt{9}$ là: \shortans{2}\par\vspace{0.6mm}
\noindent\textbf{Câu 7 (0,25 điểm).} Nghiệm của phương trình $\sqrt{3x + 1} = 4$ là $x =$ \shortans{5}\par\vspace{0.6mm}
\noindent\textbf{Câu 8 (0,25 điểm).} Thang dài $6\text{ m}$ nghiêng $60^\circ$ với đất. Chân thang cách tường: \shortans{3\text{ m}}\par\vspace{0.6mm}
\noindent\textbf{Câu 9 (0,25 điểm).} Nghiệm của hệ $\begin{cases} x + y = 6 \\ x - y = 2 \end{cases}$ là $(x; y) =$ \shortans{(4; 2)}\par\vspace{0.6mm}
\noindent\textbf{Câu 10 (0,25 điểm).} Giá trị $\cos^2 10^\circ + \dots + \cos^2 80^\circ$ (8 số hạng) là: \shortans{4}\par\vspace{0.6mm}
\vspace{0.8mm}
\noindent\textbf{Câu 11 (2,5 điểm). Đại số tổng hợp:}
\vspace{0.6mm}

\noindent
\begin{tabularx}{\linewidth}{|X|X|}
\hline
\oBaiTap[20]{a) (1,25 điểm)}{$\begin{cases} 3x + 2y = 13 \\ 2x - y = 4 \end{cases}$} & 
\oBaiTap[20]{b) (1,25 điểm)}{Rút gọn với $x > 0, x \ne 9$:\par $P = \left(\dfrac{1}{\sqrt{x}-3} + \dfrac{1}{\sqrt{x}+3}\right) \cdot \dfrac{x-9}{\sqrt{x}}$} \\
\hline
\end{tabularx}

\newpage
% =============================================================================
% ========================= TOÁN 9 - TUẦN 07 - BUỔI 02: MÃ ĐỀ 102 =============
% ========================= MẶT SAU (TRANG 2/2) ================================
% =============================================================================

\noindent
\begin{tabularx}{\linewidth}{|X|c|}
\hline
\textbf{TOÁN 9 - TUẦN 07 - BUỔI 02} & \textbf{Mã đề: 102 - Trang 2/2} \\
\hline
\end{tabularx}

\vspace{1.2mm}

\noindent
\begin{tabularx}{\linewidth}{|X|}
\hline
\vspace{0.5mm}
\textbf{Câu 12 (2,5 điểm).} Hai ô tô khởi hành cùng một lúc từ hai địa điểm $A$ và $B$ cách nhau $120\text{ km}$, đi ngược chiều nhau và gặp nhau sau $1{,}5$ giờ. Biết vận tốc ô tô đi từ $A$ lớn hơn vận tốc ô tô đi từ $B$ là $20\text{ km/h}$. Tính vận tốc của mỗi ô tô.\par
\vspace{0.4mm}
\noindent\textit{Bài giải:}\dongke{16}
\\ \hline
\end{tabularx}

\vspace{2.0mm}

\noindent
\begin{tabularx}{\linewidth}{|X|}
\hline
\vspace{0.5mm}
\textbf{Câu 13 (2,5 điểm).} Cho tam giác $ABC$ vuông tại $A$ có $\widehat{B} = 60^\circ$ và cạnh huyền $BC = 10\text{ cm}$. Kẻ đường cao $AH$ ($H \in BC$).\par
\textbf{a) (1,5 điểm)} Tính số đo góc $C$, độ dài các cạnh $AB, AC$ và độ dài đường cao $AH$.\par
\textbf{b) (1,0 điểm)} Kẻ $HD \perp AB$ tại $D$ và $HE \perp AC$ tại $E$. Chứng minh tứ giác $ADHE$ là hình chữ nhật và suy ra $DE = AH$.\par
\vspace{0.4mm}
\noindent\textit{Bài giải:}\dongke{17}
\\ \hline
\end{tabularx}

\newpage
% =============================================================================
% ========================= TOÁN 9 - TUẦN 07 - BUỔI 02: MÃ ĐỀ 103 =============
% ========================= MẶT TRƯỚC (TRANG 1/2) ==============================
% =============================================================================

\noindent
\begin{tabularx}{\linewidth}{|X|c|}
\hline
\textbf{PHIẾU KIỂM TRA ĐỊNH KỲ TOÁN 9 - TUẦN 07 - BUỔI 02} & \textbf{Mã đề: 103 - Trang 1/2} \\
\hline
\end{tabularx}

\vspace{1.0mm}

\noindent
\begin{tabularx}{\linewidth}{|p{3.2cm}|X|}
\hline
\centering\textbf{ĐIỂM SỐ} & \centering\textbf{LỜI NHẬN XÉT CỦA THẦY CÔ} \tabularnewline
\hline
\rule{0pt}{1.8cm} & \tabularnewline
\hline
\end{tabularx}

\vspace{1.2mm}
\noindent\textbf{Câu 1 (0,25 điểm).} $\Delta ABC$ vuông tại $A$ ($\widehat{B}=30^\circ, BC=14\text{ cm}$). Cạnh $AC =$ \shortans{7\text{ cm}}\par\vspace{0.6mm}
\noindent\textbf{Câu 2 (0,25 điểm).} Cho góc nhọn $\alpha$ có $\sin \alpha = \dfrac{5}{13}$. Khi đó $\cos \alpha =$ \shortans{\dfrac{12}{13}}\par\vspace{0.6mm}
\noindent\textbf{Câu 3 (0,25 điểm).} Cho $\Delta ABC$ vuông tại $A$, đường cao $AH$. Biết $BH=1, CH=9$. Khi đó $AH =$ \shortans{3}\par\vspace{0.6mm}
\noindent\textbf{Câu 4 (0,25 điểm).} Rút gọn biểu thức $\sqrt{(1 - \sqrt{3})^2}$ được kết quả: \shortans{\sqrt{3} - 1}\par\vspace{0.6mm}
\noindent\textbf{Câu 5 (0,25 điểm).} Biểu thức $\sqrt{4x - 12}$ xác định khi: \shortans{x \ge 3}\par\vspace{0.6mm}
\noindent\textbf{Câu 6 (0,25 điểm).} Giá trị của biểu thức $\sqrt{75} : \sqrt{3} + \sqrt{25}$ là: \shortans{10}\par\vspace{0.6mm}
\noindent\textbf{Câu 7 (0,25 điểm).} Nghiệm của phương trình $\sqrt{4x + 1} = 5$ là $x =$ \shortans{6}\par\vspace{0.6mm}
\noindent\textbf{Câu 8 (0,25 điểm).} Thang dài $4\text{ m}$ nghiêng $60^\circ$ với đất. Chân thang cách tường: \shortans{2\text{ m}}\par\vspace{0.6mm}
\noindent\textbf{Câu 9 (0,25 điểm).} Nghiệm của hệ $\begin{cases} x + y = 7 \\ x - y = 1 \end{cases}$ là $(x; y) =$ \shortans{(4; 3)}\par\vspace{0.6mm}
\noindent\textbf{Câu 10 (0,25 điểm).} Giá trị $\sin^2 15^\circ + \sin^2 35^\circ + \sin^2 55^\circ + \sin^2 75^\circ$ là: \shortans{2}\par\vspace{0.6mm}
\vspace{0.8mm}
\noindent\textbf{Câu 11 (2,5 điểm). Đại số tổng hợp:}
\vspace{0.6mm}

\noindent
\begin{tabularx}{\linewidth}{|X|X|}
\hline
\oBaiTap[20]{a) (1,25 điểm)}{$\begin{cases} 2x + y = 7 \\ 3x + 2y = 12 \end{cases}$} & 
\oBaiTap[20]{b) (1,25 điểm)}{Rút gọn với $x > 0, x \ne 1$:\par $P = \left(\dfrac{1}{\sqrt{x}-1} + \dfrac{1}{\sqrt{x}+1}\right) \cdot \dfrac{x-1}{\sqrt{x}}$} \\
\hline
\end{tabularx}

\newpage
% =============================================================================
% ========================= TOÁN 9 - TUẦN 07 - BUỔI 02: MÃ ĐỀ 103 =============
% ========================= MẶT SAU (TRANG 2/2) ================================
% =============================================================================

\noindent
\begin{tabularx}{\linewidth}{|X|c|}
\hline
\textbf{TOÁN 9 - TUẦN 07 - BUỔI 02} & \textbf{Mã đề: 103 - Trang 2/2} \\
\hline
\end{tabularx}

\vspace{1.2mm}

\noindent
\begin{tabularx}{\linewidth}{|X|}
\hline
\vspace{0.5mm}
\textbf{Câu 12 (2,5 điểm).} Hai người đi xe đạp cùng khởi hành một lúc từ hai địa điểm $A$ và $B$ cách nhau $48\text{ km}$, đi ngược chiều nhau và gặp nhau sau $2$ giờ. Biết vận tốc người đi từ $A$ lớn hơn vận tốc người đi từ $B$ là $4\text{ km/h}$. Tính vận tốc của mỗi người.\par
\vspace{0.4mm}
\noindent\textit{Bài giải:}\dongke{16}
\\ \hline
\end{tabularx}

\vspace{2.0mm}

\noindent
\begin{tabularx}{\linewidth}{|X|}
\hline
\vspace{0.5mm}
\textbf{Câu 13 (2,5 điểm).} Cho tam giác $ABC$ vuông tại $A$ có $\widehat{B} = 30^\circ$ và cạnh huyền $BC = 8\text{ cm}$. Kẻ đường cao $AH$ ($H \in BC$).\par
\textbf{a) (1,5 điểm)} Tính số đo góc $C$, độ dài các cạnh $AC, AB$ và độ dài đường cao $AH$.\par
\textbf{b) (1,0 điểm)} Kẻ $HD \perp AB$ tại $D$ và $HE \perp AC$ tại $E$. Chứng minh rằng: $AD \cdot AB = AE \cdot AC$.\par
\vspace{0.4mm}
\noindent\textit{Bài giải:}\dongke{17}
\\ \hline
\end{tabularx}

\newpage
% =============================================================================
% ========================= TOÁN 9 - TUẦN 07 - BUỔI 02: MÃ ĐỀ 104 =============
% ========================= MẶT TRƯỚC (TRANG 1/2) ==============================
% =============================================================================

\noindent
\begin{tabularx}{\linewidth}{|X|c|}
\hline
\textbf{PHIẾU KIỂM TRA ĐỊNH KỲ TOÁN 9 - TUẦN 07 - BUỔI 02} & \textbf{Mã đề: 104 - Trang 1/2} \\
\hline
\end{tabularx}

\vspace{1.0mm}

\noindent
\begin{tabularx}{\linewidth}{|p{3.2cm}|X|}
\hline
\centering\textbf{ĐIỂM SỐ} & \centering\textbf{LỜI NHẬN XÉT CỦA THẦY CÔ} \tabularnewline
\hline
\rule{0pt}{1.8cm} & \tabularnewline
\hline
\end{tabularx}

\vspace{1.2mm}
\noindent\textbf{Câu 1 (0,25 điểm).} $\Delta ABC$ vuông tại $A$ ($\widehat{B}=30^\circ, BC=18\text{ cm}$). Cạnh $AC =$ \shortans{9\text{ cm}}\par\vspace{0.6mm}
\noindent\textbf{Câu 2 (0,25 điểm).} Cho góc nhọn $\alpha$ có $\cos \alpha = \dfrac{5}{13}$. Khi đó $\sin \alpha =$ \shortans{\dfrac{12}{13}}\par\vspace{0.6mm}
\noindent\textbf{Câu 3 (0,25 điểm).} Cho $\Delta ABC$ vuông tại $A$, đường cao $AH$. Biết $BH=3, CH=12$. Khi đó $AH =$ \shortans{6}\par\vspace{0.6mm}
\noindent\textbf{Câu 4 (0,25 điểm).} Rút gọn biểu thức $\sqrt{(2 - \sqrt{7})^2}$ được kết quả: \shortans{\sqrt{7} - 2}\par\vspace{0.6mm}
\noindent\textbf{Câu 5 (0,25 điểm).} Biểu thức $\sqrt{5x - 15}$ xác định khi: \shortans{x \ge 3}\par\vspace{0.6mm}
\noindent\textbf{Câu 6 (0,25 điểm).} Giá trị của biểu thức $\sqrt{48} : \sqrt{3} - \sqrt{4}$ là: \shortans{2}\par\vspace{0.6mm}
\noindent\textbf{Câu 7 (0,25 điểm).} Nghiệm của phương trình $\sqrt{2x + 5} = 5$ là $x =$ \shortans{10}\par\vspace{0.6mm}
\noindent\textbf{Câu 8 (0,25 điểm).} Thang dài $8\text{ m}$ nghiêng $60^\circ$ với đất. Chân thang cách tường: \shortans{4\text{ m}}\par\vspace{0.6mm}
\noindent\textbf{Câu 9 (0,25 điểm).} Nghiệm của hệ $\begin{cases} x + y = 8 \\ x - y = 2 \end{cases}$ là $(x; y) =$ \shortans{(5; 3)}\par\vspace{0.6mm}
\noindent\textbf{Câu 10 (0,25 điểm).} Giá trị $\cos^2 15^\circ + \cos^2 35^\circ + \cos^2 55^\circ + \cos^2 75^\circ$ là: \shortans{2}\par\vspace{0.6mm}
\vspace{0.8mm}
\noindent\textbf{Câu 11 (2,5 điểm). Đại số tổng hợp:}
\vspace{0.6mm}

\noindent
\begin{tabularx}{\linewidth}{|X|X|}
\hline
\oBaiTap[20]{a) (1,25 điểm)}{$\begin{cases} 3x - 2y = 5 \\ x + 2y = 7 \end{cases}$} & 
\oBaiTap[20]{b) (1,25 điểm)}{Rút gọn với $x > 0, x \ne 16$:\par $P = \left(\dfrac{1}{\sqrt{x}-4} + \dfrac{1}{\sqrt{x}+4}\right) \cdot \dfrac{x-16}{\sqrt{x}}$} \\
\hline
\end{tabularx}

\newpage
% =============================================================================
% ========================= TOÁN 9 - TUẦN 07 - BUỔI 02: MÃ ĐỀ 104 =============
% ========================= MẶT SAU (TRANG 2/2) ================================
% =============================================================================

\noindent
\begin{tabularx}{\linewidth}{|X|c|}
\hline
\textbf{TOÁN 9 - TUẦN 07 - BUỔI 02} & \textbf{Mã đề: 104 - Trang 2/2} \\
\hline
\end{tabularx}

\vspace{1.2mm}

\noindent
\begin{tabularx}{\linewidth}{|X|}
\hline
\vspace{0.5mm}
\textbf{Câu 12 (2,5 điểm).} Hai xe máy khởi hành cùng một lúc từ hai địa điểm $A$ và $B$ cách nhau $100\text{ km}$, đi ngược chiều nhau và gặp nhau sau $1{,}25$ giờ (tức $1$ giờ $15$ phút). Biết vận tốc xe đi từ $A$ lớn hơn vận tốc xe đi từ $B$ là $16\text{ km/h}$. Tính vận tốc của mỗi xe.\par
\vspace{0.4mm}
\noindent\textit{Bài giải:}\dongke{16}
\\ \hline
\end{tabularx}

\vspace{2.0mm}

\noindent
\begin{tabularx}{\linewidth}{|X|}
\hline
\vspace{0.5mm}
\textbf{Câu 13 (2,5 điểm).} Cho tam giác $ABC$ vuông tại $A$ có $\widehat{B} = 60^\circ$ và cạnh huyền $BC = 12\text{ cm}$. Kẻ đường cao $AH$ ($H \in BC$).\par
\textbf{a) (1,5 điểm)} Tính số đo góc $C$, độ dài các cạnh $AB, AC$ và độ dài đường cao $AH$.\par
\textbf{b) (1,0 điểm)} Kẻ $HD \perp AB$ tại $D$ và $HE \perp AC$ tại $E$. Chứng minh tứ giác $ADHE$ là hình chữ nhật và suy ra $DE = AH$.\par
\vspace{0.4mm}
\noindent\textit{Bài giải:}\dongke{17}
\\ \hline
\end{tabularx}

\end{document}
